\section{The conductor-four extension in the same actual arithmetic window}
\label{sec:conductor4-four-residue-window}

\paragraph{Status and exact scope.}
This conditional-interface proof extends the established sum of
AC-G151-window using the original-coefficient strict theorem AC-G152.
The latter uses AC-G152-class only for its stated class sector and
exact field and unit data.  Their proof and calculation reviews are
recorded separately.  All unlisted coefficient primes, finite
characters and Tate weights remain in the unknown sum.
No mixed realization or bound for that remaining sum is asserted.

Put $p=269$, $K=\mathbb Q(\omega)$ and $t^3=p$ with $t$ real.
Let $\chi$ describe $K(t)/K$, and let $\phi$ be the primitive odd
character modulo sixteen with $\phi(13)=i$ and kernel $\{1,7\}$.
Write $\psi_4=\chi_{-4}|_{G_K}$ and
$\xi_\pm=\chi_{\pm8}|_{G_K}$.
Fix a source datum, $0<\varepsilon<1/256$, an allowed page and
its prescribed pair $q,d$.  Retain the actual fields and cutoff
\[
 \begin{aligned}
 F_u&=\mathbb Q(\zeta_{2^{h(u)+2}},\zeta_7,
                         \sqrt{-q},\sqrt{-d},\omega),\\
 H_u&=F_u(t),\qquad G_u=\operatorname{Gal}(H_u/K),\quad D_u=|G_u|,\\
 u^{2+8\varepsilon}&\leq2^{h(u)}\leq2u^{2+8\varepsilon},
                    \qquad T=\lfloor u^{399}\rfloor.
 \end{aligned}
 \tag{C4W.1}
\]
All original $O_\ell,E_\ell,b_\ell$, strict local conditions,
derived indices and the unrestricted coefficient-prime range are
those of (FC.1)--(FC.2), (EU.2) and (TC.1).

For $u\geq2$ use the following disjoint label sets:
\[
 \begin{aligned}
 Z_{16}^+&=\{(p,\chi^{-j}\phi^{-\sigma},e):
   j=1,2,\ \sigma=\pm1,\\
 &\hspace{2em}1\leq e<T,\ e\equiv-1\pmod{67}\},\\
 Z_4^+&=\{(p,\chi^{-j}\psi_4,e):
   j=1,2,\ 1\leq e<T,\ e\equiv-1\pmod{67}\},\\
 Z_8&=\{(p,\chi^{-j}\xi_\nu,e):
   j=1,2,\ \nu=\pm,\ 1\leq e<T,\ e\equiv-1\pmod{134}\},\\
 Z_{\rm cub}&=\{(p,\chi^j,e):j=0,1,2,
   \ 1\leq e<T,\ e\equiv0,267\pmod{268}\},\\
 Z_3&=\{(3,1,e):1\leq e<T\},\\
 Z^{4+}&=Z_{16}^+\sqcup Z_4^+\sqcup Z_8\sqcup Z_{\rm cub}\sqcup Z_3.
 \end{aligned}
 \tag{C4W.2}
\]
The former conductor-four set $Z_4$ in AC-G151 has the condition
$e\equiv-1\pmod{134}$ instead.  Thus precisely the new set is
\[
 Z_4^{\rm new}=\{(p,\chi^{-j}\psi_4,e):j=1,2,
       \ 1\leq e<T,\ e\equiv66,200\pmod{268}\}.
 \tag{C4W.3}
\]
Let $I_{399}^{\rm est4+}$ be (TC.1) with exactly $Z^{4+}$ omitted.
The old remaining sum is $I_{399}^{\rm est8}$ from AC-G151.

\begin{theorem}[AC-G153-window: the actual conductor-four extension]
\label{AC-G153-window}
Every label in (C4W.2) occurs on the prescribed host, and its
strict group on the original coefficient module vanishes.
For positive integer $T$ the exact counts are
\[
 \begin{aligned}
 |Z_{16}^+|&=4\lfloor T/67\rfloor,&
 |Z_4^+|&=2\lfloor T/67\rfloor,\\
 |Z_8|&=4\lfloor T/134\rfloor,& |Z_3|&=T-1,\\
 |Z_{\rm cub}|&=3\bigl(\lfloor(T-1)/268\rfloor
                                +\lfloor T/268\rfloor\bigr),\\
 |Z_4^{\rm new}|&=2\lfloor(T+67)/134\rfloor.
 \end{aligned}
 \tag{C4W.4}
\]
Let $C_u$ be the sum of the actual nonnegative normalized
positive-index costs over the entire union $Z^{4+}$.  Then
\[
 \begin{aligned}
 I_{399}&=I_{399}^{\rm est4+}+C_u,\qquad
 I_{399}^{\rm est4+}\leq I_{399}^{\rm est8}\leq I_{399},\\
 0\leq C_u&\leq8D_u^2T^2\log(8TD_u)
             +\tfrac12D_u\mathcal R_uT^2+2D_ur_uc_0T\\
 &=O_\varepsilon\bigl(u^{802+16\varepsilon}(\log u)^2\bigr).
 \end{aligned}
 \tag{C4W.5}
\]
The equality is an identity of nonnegative extended sums;
no subtraction of infinities is used.  One total (FL.4) majorant
covers the entire old and new union.

The common-window hypothesis (TC.3), or its AC-G151 form, may
therefore use $I_{399}^{\rm est4+}$ with all its simultaneous
field, prime, character and Tate-weight quantifiers unchanged.
Either old hypothesis implies the new one with the same parameters.
Conversely the new hypothesis implies the full old one using the
same $W,J$, source datum, page and pair, after increasing only
the starting threshold and halving its positive relative slack.
The actual source $R.t$, the mixed theta/Selmer and analytic input
$\mathcal R_{\rm cub}$, and every other AC-K35 hypothesis remain.
None of these unproved inputs is supplied by this theorem.
\end{theorem}

\begin{proof}
\proofstep{1. The actual finite characters and the original coefficients.}
For $u\geq2$, (C4W.1) gives $h(u)\geq2$, hence
$\mathbb Q(\zeta_{16})\subset F_u$.
In particular $\psi_4$ factors through the same $G_u$;
there is no change of host.  Since $F_u/\mathbb Q$ is abelian,
it cannot contain the nonabelian $S_3$ field $K(t)$.
As $K(t)/K$ has prime degree three, $F_u\cap K(t)=K$.
Its cubic subgroup distinguishes $j=1,2$.
The trivial, conductor-four and two conductor-eight quadratic
characters stay distinct on $G_K$: any equality would make a
nontrivial quadratic character ramified only at two equal the
character of $K/\mathbb Q$, which ramifies at three.
The conductor-sixteen characters have order four and are distinct
from all these quadratic factors.  The coefficient-three set
has a different prime.  Thus all five sets in (C4W.2) are disjoint;
their old occurrence assertions are exactly those of AC-G151.

The original $O_p$ contains the cubic values, as
$\exp(G_u)=3\cdot2^{h(u)}$, and $\psi_4$ has rational values.
For a new label, its coefficient module is literally
\[
 \begin{aligned}
 A_{j,e}&=(O_p\otimes_{\mathbb Z_p}\mathbb Q_p/\mathbb Z_p)
                         (\chi^{-j}\psi_4)(-e)\\
        &=(E_p/O_p)(\chi^{-j}\psi_4c_p^{-e}).
 \end{aligned}
 \tag{C4W.6}
\]
Apply AC-G152 with its additional exponent $m=0$ and
continuous character $\kappa=1$.
It permits this original $O_p$, with arbitrary absolute
ramification, and gives strict zero for $e\equiv66,200\bmod268$.
Its conductor-$269$ quartic character describes only the residual
Kummer cover.  Neither that character nor its quadratic square
is adjoined to (C4W.6) or to the finite host labels.
The old residues $133,267\bmod268$ are retained from AC-G151.
Every other old strict assertion, including the coefficient-three
comparison without division by three, remains unchanged.

\proofstep{2. Full local conditions and exact counts.}
AC-G152 imposes zero in full $H^1(G_{K_v},A_{j,e})$ at $v\mid p$
and zero in $H^1(I_v,A_{j,e})$ at every other finite place.
These are exactly the strict conditions of (FC.2).
The coefficient (C4W.6) is unramified away from $2,3,p$.
At any added away-$p$ prime in $\Sigma_p$, inertia acts trivially,
so a strict restriction that is a coboundary is actually zero.
The cocycle identity makes the restriction zero on all conjugate
inertia groups and the closed normal subgroup they generate.
The action and cocycle therefore descend to the smaller finite
ramification set.  Conversely inflation is injective and preserves
the other local conditions.  This identifies the strict kernels
without replacing their full local condition at $p$.

The two new residues are $e=66+134n$, $n\geq0$.
For integer $T>0$ the number of these with $1\leq e<T$ is
$\lfloor(T+67)/134\rfloor$.
Their union with $133,267\bmod268$ is exactly $e=67n-1$,
$n\geq1$, whose count below $T$ is $\lfloor T/67\rfloor$.
Multiplication by the two distinct cubic labels gives (C4W.4).
In particular
\[
 \lfloor T/67\rfloor=\lfloor T/134\rfloor
                          +\lfloor(T+67)/134\rfloor.
 \tag{C4W.7}
\]
All other counts are the unchanged AC-G151 counts.
Thus precisely the indicated strict terms in (FC.15) vanish.
Strict zero does not assert that a derived positive index is zero.

\proofstep{3. One full-union local comparison cost.}
For any omitted $(\ell,\eta,e)$ use the original dual complex
$C_{\ell,\eta,e}=R\Gamma(G_{K,\Sigma_\ell},
O_\ell(\eta^{-1})(e+1))$ and exactly the codifferent convention
of (FL.2).  The determinant inequality (EU.2) gives
\[
 \begin{aligned}
 \frac{\log\#k_{O_\ell}}{b_\ell}\max(0,i_{\ell,\eta,e})
 &\leq\frac{\log\#H^2(C_{\ell,\eta,e})}{b_\ell}\\
 &=\frac{|\log\#H^2(C_{\ell,\eta,e})
                    -\log\#\mathcal S_{\ell,\eta}(e)|}{b_\ell}.
 \end{aligned}
 \tag{C4W.8}
\]
The equality uses the proved strict zero only.
At a fixed weight, the sum over the entire union in (C4W.2)
is a single subsum of the nonnegative all-prime error (FL.4),
including its coefficient-three terms.  It is bounded by
\[
       16eD_u^2\log(8eD_u)+D_ue\mathcal R_u+2D_ur_uc_0.
 \tag{C4W.9}
\]
Summing once over $1\leq e<T$, using $\sum e\leq T^2/2$,
$\log(8eD_u)\leq\log(8TD_u)$ and $T-1\leq T$, proves (C4W.5).
There is no additional full comparison cost for each family.
Nonnegativity proves the additive identity and monotonicity even
if one of the remaining sums is infinite.

The same host bounds give $D_u\leq288u^{2+8\varepsilon}$,
$r_u\leq12$, and
$\mathcal R_u=O_\varepsilon(u^{\varepsilon/5}(\log u)^2)$,
uniformly over the prescribed page and pair.
The two largest exponents after $T\leq u^{399}$ are
$802+16\varepsilon$ and $800+41\varepsilon/5$.
Thus the error is $o(\exp(u^\varepsilon/2))$ for each fixed
positive $\varepsilon$, with the same uniformity.

\proofstep{4. Preserve the common window and every realization input.}
Fix the datum, page and $0<\varepsilon<1/256$.
Suppose the new premise supplies $J\geq1$, $X_0$ and
$s\in(0,1)$, independent of $X,q,d,u,\ell,\eta,e$, such that
for every $X\geq X_0$ and prescribed
$\operatorname{WinPair}(\varepsilon,X^{1/J},q,d)$ there is one
$W$ with $X^{1/J}\leq W$, $W^2\leq X$ and
\[
 I_{399}^{\rm est4+}(H_u,u)\leq(1-s)\exp(u^\varepsilon/2)
                                  \quad(W\leq u\leq W^2).
 \tag{C4W.10}
\]
Choose $U_0\geq2$ so that (C4W.5) is at most
$(s/2)\exp(u^\varepsilon/2)$ for all $u\geq U_0$.
It depends only on the retained datum, page, $\varepsilon$ and
slack, not on a new pair, coefficient, character, weight or window.
Increasing $X_0$ to $\max(X_0,U_0^J)$ keeps the very same $W$
and gives
\[
 I_{399}(H_u,u)\leq(1-s/2)\exp(u^\varepsilon/2)
                                  \quad(W\leq u\leq W^2).
 \tag{C4W.11}
\]
For a fixed finite simultaneous family take the maximum of its
error constants, retaining one window; no intersection of separately
chosen windows is invoked.  Conversely either the full old premise
or AC-G151's premise implies (C4W.10) by nonnegativity, with its
original parameters.  Uniform little-$o$ formulations agree too.
Nothing changes the actual source host, rank $R.t$, family
construction, mixed theta/Selmer identification or analytic
realization clause in $\mathcal R_{\rm cub}$.
\end{proof}

\paragraph{Precise remaining gate and geometric scope.}
The strongest old EB-Q02 premise after AC-G151 still includes
$Z_4^{\rm new}$.  Its count is positive for $T\geq67$,
so the unknown index sum omits a proper additional subset of
actual terms.  This does not assert a strict numerical inequality
between the two sums.  Every old omission is retained without
being counted again as new progress.
The remaining all-prime bound, mixed realization, analytic
comparison and H04.4 remain open.  All coefficients are the
original rank-one Artin--Tate systems on the same arithmetic
curve and finite cover.  The previously audited SGA~1,
Expos\'e~V, \S\S4 and 7 dictionary is retained only at that scope;
no new motive, comparison theorem or Lean execution is claimed.
